\question[4]
Gegeben ist die folgende Tabelle einer Schule:

\begin{center}
\begin{tabular}{|c|l|l|l|}
\hline
SchuelerID & Schuelername & Klassenname & Klassenlehrer \\
\hline
1 & Anna Meier & 10A & Herr Müller \\
2 & Max Schulz & 10A & Herr Müller \\
3 & Lea Weber & 10B & Frau Schmidt \\
\hline
\end{tabular}
\end{center}

Erläutere, warum die Tabelle Redundanzen enthält.
Überführe die Tabelle anschließend in ein relationales Datenbankschema in der 3. Normalform.

\begin{solutionbox}{5cm}
Die Information, wer der Klassenlehrer eine Klasse ist, wird mehrfach gespeichert. Die Information 
„10A -- Herr Müller“ erscheint in mehreren Datensätzen.
Dadurch können Redundanzen und Inkonsistenzen entstehen.

\vspace{0.3cm}

\textbf{Relationales Schema:}

\begin{itemize}
\item Schueler(\pk{SchuelerID}, Schuelername, \fk{KlassenID})
\item Klasse(\pk{KlassenID}, Klassenname, Klassenlehrer)
\end{itemize}

Die Informationen über Klassen werden nun nur noch einmal gespeichert.
\end{solutionbox}